(1) 往返时间 $t_{\rm rt}=2L/c$。一次往返穿过介质距离 $2L$，能量存留因子 $r=(1-T_1)(1-T_2)e^{-2\alpha L}$；假定镜面额外散射已并入 $T_i$，并且每次往返的损耗相同。
(2) 往返 $n$ 次后 $U_n=U_0r^n$，以 $t=nt_{\rm rt}$ 连续化得 $U(t)=U_0e^{-t/\tau}$，其中 $\tau=-t_{\rm rt}/\ln r$。当 $T_1,T_2,\alpha L\ll1$，$\tau\simeq2L/[c(T_1+T_2+2\alpha L)]$。
(3) 定义 $Q=\omega_0U/P_{\rm loss}$，由于 $P_{\rm loss}=U/\tau$，有 $Q=\omega_0\tau$。场振幅衰减为 $e^{-t/(2\tau)}$，Lorentz共振强度的频率半高全宽 $\Delta\nu=1/(2\pi\tau)$；不能把振幅寿命和能量寿命混用。
(4) 镜损耗与腔长相同、仅改变吸收时，$\tau_s^{-1}-\tau_0^{-1}=c\alpha$，故 $\alpha=(\tau_s^{-1}-\tau_0^{-1})/c$。该差分在弱吸收下消去固定镜损耗；若插入样品改变腔长、反射或散射，必须额外校准。
